\chapter{Wstęp}
\label{chap:wstep}

Komunikatory głosowe stały się powszechnym miejscem spotkań zarówno zespołów
projektowych, jak i~społeczności skupionych wokół wspólnych zainteresowań.
Jedną z~najpopularniejszych platform tego typu jest Discord~\cite{Discord},
który pozwala prowadzić rozmowy na kanałach głosowych serwerów. Platforma nie
oferuje jednak wbudowanego mechanizmu tworzenia transkrypcji ani notatek ze
spotkań. Ustalenia, które zapadają w~trakcie rozmowy, pozostają więc wyłącznie
w~pamięci uczestników albo w~notatkach sporządzanych ręcznie.

Problem ten jest szczególnie widoczny podczas zdalnie prowadzonych sesji gier
fabularnych (ang.~\textit{role-playing games}, \acs{rpg}), takich jak
\textit{Dungeons \& Dragons}. Sesja trwa zazwyczaj kilka godzin, a~w~jej trakcie
pojawia się wiele informacji, które trzeba pamiętać na kolejnych spotkaniach:
imiona postaci niezależnych, odwiedzone miejsca, zdobyte przedmioty, podjęte
decyzje oraz niezakończone wątki fabularne. Ręczne notowanie odciąga gracza od
rozgrywki, a~powstałe w~ten sposób notatki są zwykle niepełne.

Rozwój metod sztucznej inteligencji pozwala ten proces zautomatyzować. Modele
automatycznego rozpoznawania mowy, takie jak Whisper~\cite{Whisper}, obsługują
wiele języków, w~tym język polski, a~duże modele językowe (ang.~\textit{Large
Language Models}, \acs{llm}) potrafią na podstawie długiego i~niedoskonałego
tekstu przygotować zwięzłe, uporządkowane streszczenie. Oba rodzaje modeli
wymagają jednak znacznej mocy obliczeniowej, najlepiej procesora graficznego
(ang.~\textit{Graphics Processing Unit}, \acs{gpu}). Utrzymywanie stale
włączonej maszyny z~\acs{gpu} na potrzeby spotkań, które odbywają
się nieregularnie, byłoby nieopłacalne. Istotną częścią problemu jest więc także
dobór infrastruktury chmurowej, w~której zasoby obliczeniowe są uruchamiane
tylko wtedy, gdy są potrzebne, a~płaci się wyłącznie za faktyczny czas ich
pracy.

\section{Cel i zakres pracy}
\label{sec:wstep:cel}

Celem pracy jest zaprojektowanie i~implementacja rozproszonego systemu, który
będzie automatycznie tworzył notatki ze spotkań prowadzonych na kanałach
głosowych platformy Discord. System będzie nagrywał rozmowę, tworzył jej
transkrypcję z~podziałem na mówców, generował na jej podstawie uporządkowane
notatki w~formacie Markdown i~publikował je na kanale tekstowym serwera.
Użytkownik będzie obsługiwał system wyłącznie za pomocą komend wpisywanych na
kanale tekstowym Discorda. Jako główny przypadek użycia zostaną przyjęte sesje
gry \textit{Dungeons \& Dragons} prowadzone w~języku polskim.

Realizacja celu obejmuje utworzenie systemu działającego w~rozproszonej
infrastrukturze chmury Microsoft Azure. Na tej platformie zostaną utworzone
niezbędne zasoby, dobrane z~uwzględnieniem kosztów utrzymania oraz czasu
działania. Poza samym połączeniem z~Discordem, z~którego będą pochodzić
nagrania, całe przetwarzanie danych będzie się odbywać w~obrębie chmury Azure,
a~poszczególne komponenty systemu będą się komunikować za pośrednictwem jej
usług. Założono podział systemu na dwie części. Pierwsza, o~niskim zapotrzebowaniu
na zasoby, będzie odpowiadać za komunikację z~Discordem i~będzie działać stale
na taniej maszynie wirtualnej. Druga, o~dużej mocy obliczeniowej, będzie
wykonywać kosztowne obliczenia z~użyciem modeli sztucznej inteligencji
i~będzie uruchamiana tylko na żądanie. Obie części będą wymieniać dane
wyłącznie przez usługę Azure Blob Storage.

Przetwarzanie nagrań zostanie powierzone agentom sztucznej inteligencji
(ang.~\textit{Artificial Intelligence}, \acs{ai}), z~których każdy będzie miał
ściśle określoną rolę, dane wejściowe i~oczekiwany wynik. Pierwszy z~nich,
oparty na modelu automatycznego rozpoznawania mowy (ang.~\textit{Automatic
Speech Recognition}, \acs{asr}) Whisper, będzie zamieniał nagrania
poszczególnych uczestników na tekst. Kolejne dwa będą wykorzystywały ten sam,
lokalnie uruchamiany model \acs{llm}, ale z~odmiennymi instrukcjami
systemowymi. Agent streszczający będzie analizował
kolejne fragmenty transkrypcji i~wyciągał z~nich najważniejsze wydarzenia,
decyzje, postacie i~miejsca, pomijając błędy rozpoznawania mowy.
Agent redagujący będzie łączył te streszczenia w~jeden spójny dokument
o~stałej strukturze: podsumowanie sesji, kluczowe informacje na przyszłość
oraz otwarte wątki. Podział zadania na mniejsze etapy
pozwoli przetwarzać nagrania trwające kilka godzin, których transkrypcja nie
zmieściłaby się w~oknie kontekstowym modelu \acs{llm}.

W~ramach pracy zaplanowano realizację następujących celów szczegółowych:
\begin{itemize}
  \item implementacja bota w~środowisku Node.js z~wykorzystaniem bibliotek
        \texttt{discord.js} i~\texttt{@discordjs/voice}, który na komendę
        użytkownika będzie dołączał do kanału głosowego i~nagrywał osobny
        strumień audio dla każdego uczestnika, dzieląc nagrania na fragmenty po
        wykryciu ciszy,
  \item zaprojektowanie sposobu przechowywania nagrań i~wyników w~usłudze
        Azure Blob Storage oraz mechanizmu śledzenia stanu każdej sesji
        (od nagrania, przez transkrypcję, po gotowe notatki lub błąd),
  \item implementacja modułu transkrypcji wykorzystującego model
        Whisper \texttt{large-v3}, który będzie łączył wypowiedzi wszystkich
        uczestników w~jeden chronologiczny dialog z~oznaczeniem mówców,
  \item implementacja modułu generowania notatek z~użyciem lokalnie
        uruchamianego modelu językowego Gemma~\cite{Gemma} udostępnianego przez
        serwer Ollama~\cite{Ollama}; moduł będzie dzielił transkrypcję na
        fragmenty, streszczał każdy z~nich, a~następnie scalał streszczenia
        w~notatki o~stałej strukturze,
  \item wdrożenie obu modułów w~chmurze Azure na zasobach z~dostępem do
        \acs{gpu}, uruchamianych wyłącznie na czas przetwarzania nagrań,
  \item integracja bota z~częścią przetwarzającą, obejmująca automatyczne
        uruchamianie przetwarzania po zakończeniu nagrania, monitorowanie jego
        postępu i~publikowanie gotowych notatek na kanale tekstowym,
  \item ograniczenie kosztów działania systemu, m.in. przez korzystanie z~\acs{gpu}
        tylko w~czasie przetwarzania i~usuwanie surowych nagrań po jego
        zakończeniu,
  \item przeprowadzenie testów systemu oraz ocena jakości transkrypcji
        i~notatek, a~także czasu i~kosztu przetwarzania sesji.
\end{itemize}

Zakres pracy nie obejmuje trenowania ani dostrajania modeli sztucznej
inteligencji. W~systemie zostaną wykorzystane gotowe, publicznie dostępne
modele, a~ich zachowanie zostanie dostosowane do zadania za pomocą odpowiednio
przygotowanych instrukcji (ang.~\textit{prompts}) i~parametrów generowania.

\section{Układ pracy}
\label{sec:wstep:uklad}

% UWAGA: opis rozdziałów odpowiada układowi z wzoru promotora
% (Wprowadzenie, Implementacja, Testy, Podsumowanie). Popraw, jeśli
% struktura pracy się zmieni.

Praca składa się z~pięciu rozdziałów. W~rozdziale pierwszym przedstawiono
tematykę pracy, uzasadnienie wyboru zagadnienia oraz cel i~zakres pracy.

Rozdział drugi stanowi wprowadzenie do zagadnień, na których opiera się
system. Omówiono w~nim istniejące narzędzia do transkrypcji i~notowania
spotkań, rozwiązania pozwalające nagrywać rozmowy na platformie Discord oraz
wykorzystane technologie: modele rozpoznawania mowy, duże modele językowe
i~usługi chmury Microsoft Azure.

W~rozdziale trzecim opisano implementację systemu. Przedstawiono jego
architekturę, działanie bota nagrywającego rozmowy, sposób przechowywania
danych i~śledzenia stanu sesji, moduły transkrypcji i~generowania notatek
oraz konteneryzację i~wdrożenie w~chmurze. Omówiono również pierwotny wariant
wdrożenia oparty na dwóch maszynach wirtualnych oraz powody przejścia na
zadanie uruchamiane na bezserwerowym \acs{gpu}.

Rozdział czwarty poświęcono testom systemu. Opisano w~nim sposób testowania
oraz wyniki dotyczące jakości transkrypcji i~notatek, czasu przetwarzania
i~kosztów działania.

Rozdział piąty zawiera podsumowanie pracy, ocenę stopnia realizacji celu,
ograniczenia przyjętego rozwiązania oraz możliwe kierunki jego dalszego
rozwoju.

\section{Wykorzystanie AI}
\label{sec:wstep:ai}

% Sekcja obowiązkowa wg szablonu WIT -- zob. wytyczne \cite{WITWytyczneSI}.
% TODO: uzupełnij, czy i jak korzystałeś z narzędzi AI przy pisaniu kodu,
% oraz oznacz takie fragmenty w komentarzach w kodzie.

Podczas przygotowania pracy, zgodnie z~wytycznymi
Wydziału~\cite{WITWytyczneSI}, korzystano z~asystenta Claude Opus~5.5 firmy
Anthropic (w~środowisku Claude Code). Asystent posłużył do analizy
repozytorium projektu oraz do przygotowania wstępnych wersji streszczenia
i~wstępu, które następnie zostały zredagowane przez autora.

Należy też podkreślić, że modele sztucznej inteligencji (Whisper oraz Gemma)
są częścią samego systemu będącego przedmiotem pracy. Ich zastosowanie
opisano w~rozdziale trzecim.
